\providecommand{\Ric}{\operatorname{Ric}}
\providecommand{\divergence}{\operatorname{div}}
\providecommand{\grad}{\operatorname{grad}}
\providecommand{\Exp}{\operatorname{exp}}
\providecommand{\Log}{\operatorname{log}}
\providecommand{\TpM}{T_p\mathcal{M}}
\providecommand{\M}{\mathcal{M}}
\providecommand{\N}{\mathbb{N}}

\chapter{Geod\'esicas: Aprendizaje en Variedades}
\label{ch:geodesics}

Los cap\'itulos anteriores construyeron dos pilares del aprendizaje profundo
geom\'etrico: los \emph{dominios} donde residen los datos y los \emph{grafos}
que los discretizan. El \Cref{ch:graphs}, en particular, mostr\'o que el
laplaciano de un grafo de proximidad converge al operador de Laplace--Beltrami
de la superficie subyacente, y aprovech\'o esta conexi\'on para definir
transformadas de Fourier intr\'insecas y filtros espectrales. Ese operador, sin
embargo, presupone que el dominio est\'a dotado de una \emph{m\'etrica
riemanniana}, noci\'on que las variedades diferenciables del
\Cref{ch:foundations} a\'un no poseen: tienen estructura topol\'ogica y
diferenciable, pero carecen de un sentido intr\'inseco de distancia.

Las \emph{geod\'esicas} completan este vac\'io. En el espacio euclidiano, el
camino m\'as corto entre dos puntos es trivialmente una recta; sobre una
variedad curva, las geod\'esicas generalizan esta idea como curvas que minimizan
localmente la longitud. Son el hilo que une los temas precedentes: proporcionan
las distancias intr\'insecas que justifican los pesos de los grafos de
proximidad, definen los sistemas de coordenadas locales necesarios para las
convoluciones geom\'etricas y gobiernan c\'omo se propaga la informaci\'on de
forma compatible con la curvatura del dominio. Para precisar estas ideas,
dotamos primero a las variedades diferenciables de estructura m\'etrica
(\Cref{sec:geo:intrinsic}), introducimos la conexi\'on y la curvatura que miden
la desviaci\'on respecto al caso plano, definimos las geod\'esicas junto con los
mapas exponencial y logar\'itmico (\Cref{sec:geo:geodesics}) y luego pasamos al
aprendizaje: la convoluci\'on geod\'esica (\Cref{sec:geo:convolution}), la
discretizaci\'on de las geod\'esicas en grafos (\Cref{sec:geo:graphs}) y los
m\'etodos espectrales en variedades (\Cref{sec:geo:spectral}).

\section{Geometr\'ia intr\'inseca}
\label{sec:geo:intrinsic}

\subsection{M\'etricas riemannianas}
\label{subsec:geo:metrics}

Para hablar de geod\'esicas con rigor debemos equipar las variedades
diferenciables del \Cref{ch:foundations} con estructura m\'etrica. Una
\emph{m\'etrica riemanniana} proporciona la maquinaria para medir longitudes,
\'angulos y vol\'umenes de manera intr\'inseca, es decir, usando solo magnitudes
disponibles en la propia variedad, sin referencia a un espacio ambiente.

\begin{definition}[M\'etrica riemanniana]
\label{def:geo:metric}
Una \emph{m\'etrica riemanniana} en una variedad diferenciable $\M$ es una
asignaci\'on suave que a cada punto $p \in \M$ le asocia un producto interior
$g_p \colon \TpM \times \TpM \to \R$ en el espacio tangente. En coordenadas
locales $(x^1, \dots, x^m)$ se escribe
\[
  g_p = \sum_{i,j=1}^m g_{ij}(p)\, dx^i \otimes dx^j ,
\]
donde las $g_{ij}(p)$ son funciones suaves que forman una matriz sim\'etrica
definida positiva en cada punto. El par $(\M, g)$ es una \emph{variedad
riemanniana}.
\end{definition}

\begin{example}[M\'etricas en variedades comunes]
\label{ex:geo:metrics}
\leavevmode
\begin{itemize}
  \item \textbf{Espacio euclidiano} $\R^n$: la m\'etrica est\'andar es
    $g_{ij} = \delta_{ij}$, lo que da $ds^2 = dx_1^2 + \cdots + dx_n^2$.
  \item \textbf{Esfera} $S^2$: en coordenadas esf\'ericas $(\theta, \phi)$, la
    m\'etrica inducida por la inmersi\'on en $\R^3$ es
    $ds^2 = d\theta^2 + \sin^2\theta\, d\phi^2$.
  \item \textbf{Espacio hiperb\'olico} $\mathbb{H}^n$ (modelo de Poincar\'e): en
    el disco $\mathbb{D}^n = \{x \in \R^n : \norm{x} < 1\}$ la m\'etrica es
    $g_{ij} = 4\delta_{ij} / (1 - \norm{x}^2)^2$, que ``estira'' las distancias
    cerca del borde.
\end{itemize}
\end{example}

La m\'etrica proporciona todo lo necesario para medir longitud y distancia,
mediante una cadena de tres pasos. Primero, como $g_p$ es un producto interior
en cada espacio tangente, induce una norma puntual $\norm{v}_g = \sqrt{g_p(v,v)}$;
para una curva $\gamma(t)$, esta mide la velocidad instant\'anea
$\norm{\dot\gamma(t)}_g$. Segundo, integrando la velocidad a lo largo de la
curva se obtiene su longitud total. Tercero, tomando el \'infimo de la longitud
sobre todas las curvas que unen dos puntos se obtiene la distancia geod\'esica.

\begin{definition}[Longitud y distancia geod\'esica]
\label{def:geo:distance}
La \emph{longitud} de una curva $\gamma \colon [a,b] \to \M$ es
\[
  L(\gamma) = \int_a^b \sqrt{g_{\gamma(t)}\big(\dot\gamma(t), \dot\gamma(t)\big)}\, dt ,
\]
y la \emph{distancia geod\'esica} entre $p, q \in \M$ es
\[
  d_g(p,q) = \inf\{ L(\gamma) : \gamma \text{ curva suave de } p \text{ a } q\} .
\]
\end{definition}

\begin{proposition}[La distancia geod\'esica es una m\'etrica]
\label{prop:geo:dg-metric}
En una variedad riemanniana conexa $(\M, g)$, la funci\'on $d_g$ es una
m\'etrica en sentido topol\'ogico: para todo $p, q, r \in \M$,
$d_g(p,q) \ge 0$ con igualdad si y solo si $p = q$; $d_g(p,q) = d_g(q,p)$; y
$d_g(p,r) \le d_g(p,q) + d_g(q,r)$.
\end{proposition}

La simetr\'ia se sigue de invertir la parametrizaci\'on de las curvas, y la
desigualdad triangular de concatenar curvas casi \'optimas. La no degeneraci\'on
es el \'unico punto delicado: si $p \neq q$, un argumento de compacidad en una
bola coordenada muestra que $g_x(v,v) \ge c\,\norm{v}_{\mathrm{euc}}^2$ para
alg\'un $c > 0$, de modo que toda curva que abandone la bola tiene longitud
acotada inferiormente por una constante positiva independiente de la curva.

\begin{example}[Distancia geod\'esica en la esfera]
\label{ex:geo:sphere-distance}
En la esfera unitaria $S^2$ con m\'etrica $ds^2 = d\theta^2 + \sin^2\theta\,
d\phi^2$, el arco de gran c\'irculo del polo norte $p$ a
$q = (\sin\alpha, 0, \cos\alpha)$ es $\gamma(t) = (\sin t, 0, \cos t)$,
$t \in [0,\alpha]$. Su velocidad es $g_{\gamma(t)}(\dot\gamma, \dot\gamma) = 1$,
luego $L(\gamma) = \alpha$. Como los arcos de gran c\'irculo minimizan la
longitud, $d_g(p,q) = \alpha = \arccos\inner{p}{q}$, donde
$\inner{\cdot}{\cdot}$ es el producto interior euclidiano de $\R^3$.
\end{example}

La m\'etrica tambi\'en da \'angulos, $\cos\theta = g_p(v,w) /
(\norm{v}_g \norm{w}_g)$, y un \emph{elemento de volumen} intr\'inseco
$dV_g = \sqrt{\det(g_{ij})}\, dx^1 \wedge \cdots \wedge dx^m$, que permite
integrar funciones sobre la variedad.

\subsection{El operador de Laplace--Beltrami}
\label{subsec:geo:laplacian}

El laplaciano euclidiano $\nabla^2 f = \sum_i \partial^2 f / \partial(x^i)^2 =
\divergence(\grad f)$ mide cu\'anto se desv\'ia el valor de $f$ en un punto
respecto de su promedio local. Reemplazando el gradiente y la divergencia
euclidianos por sus an\'alogos m\'etricos se obtiene su generalizaci\'on a
variedades.

\begin{definition}[Operador de Laplace--Beltrami]
\label{def:geo:laplace-beltrami}
En una variedad riemanniana $(\M, g)$, el \emph{operador de Laplace--Beltrami}
$\Delta_\M$ act\'ua sobre funciones suaves $f \colon \M \to \R$ como
\[
  \Delta_\M f = \divergence(\grad f)
    = \frac{1}{\sqrt{\det g}}\, \partial_i\!\left(\sqrt{\det g}\; g^{ij}\, \partial_j f\right),
\]
donde $g^{ij}$ son las componentes de la m\'etrica inversa. Cuando
$(\M, g) = (\R^m, \delta_{ij})$ se tiene $\det g = 1$ y $g^{ij} = \delta^{ij}$,
de modo que $\Delta_\M$ se reduce al laplaciano euclidiano.
\end{definition}

El operador de Laplace--Beltrami es el \'unico objeto que enlaza la geometr\'ia
diferencial con el aprendizaje profundo geom\'etrico, por cuatro razones.

\begin{remark}[Por qu\'e importa el operador de Laplace--Beltrami]
\label{rem:geo:lb-role}
\leavevmode
\begin{enumerate}
  \item \textbf{Difusi\'on.} $\Delta_\M$ gobierna la ecuaci\'on del calor
    $\partial_t u = \Delta_\M u$, que modela la propagaci\'on de informaci\'on
    por la variedad. Una capa que ``difunde'' caracter\'isticas hacia
    vecindarios geod\'esicos es el an\'alogo aprendible de este proceso.
  \item \textbf{An\'alisis espectral intr\'inseco.} Las autofunciones
    $\{\varphi_k\}$ de $\Delta_\M$, con $\Delta_\M \varphi_k = \lambda_k
    \varphi_k$, forman una base ortonormal de $L^2(\M)$ que generaliza la base
    de Fourier, lo que habilita una transformada de Fourier intr\'inseca
    $\hat f(\lambda_k) = \inner{f}{\varphi_k}$ y, con ella, convoluciones
    espectrales (\Cref{sec:geo:spectral}).
  \item \textbf{Invarianza bajo isometr\'ias.} Para toda isometr\'ia $\phi$ de
    $\M$, $\Delta_\M(f \circ \phi) = (\Delta_\M f) \circ \phi$. Cualquier
    arquitectura construida a partir de $\Delta_\M$ hereda as\'i la invarianza
    bajo las simetr\'ias de la variedad, justo el principio de equivarianza del
    \Cref{ch:priors}.
  \item \textbf{Discretizaci\'on.} El laplaciano de grafo del \Cref{ch:graphs}
    es el an\'alogo discreto de $\Delta_\M$ y, bajo muestreo adecuado, converge
    a \'el, de modo que las t\'ecnicas se transfieren entre los mundos continuo
    y discreto.
\end{enumerate}
\end{remark}

\subsection{Conexiones y curvatura}
\label{subsec:geo:curvature}

Para definir geod\'esicas y cuantificar cu\'anto se desv\'ia una variedad del
espacio plano necesitamos dos herramientas m\'as. En $\R^n$ todos los espacios
tangentes se identifican can\'onicamente con $\R^n$, de modo que los vectores en
puntos distintos pueden compararse directamente; en una variedad general $\TpM$
y $T_q\M$ son espacios vectoriales distintos sin identificaci\'on can\'onica.
Una \emph{conexi\'on} repara esto al prescribir c\'omo derivar campos
vectoriales, y el \emph{tensor de curvatura} mide entonces la no conmutatividad
de las derivadas covariantes resultantes.

\begin{definition}[Conexi\'on af\'in]
\label{def:geo:connection}
Una \emph{conexi\'on af\'in} en $\M$ es una aplicaci\'on
$\nabla \colon \Gamma(T\M) \times \Gamma(T\M) \to \Gamma(T\M)$,
$(X,Y) \mapsto \nabla_X Y$ sobre campos vectoriales suaves que es
$C^\infty(\M)$-lineal en $X$, $\R$-lineal en $Y$, y satisface la regla de
Leibniz $\nabla_X(fY) = (Xf)Y + f\nabla_X Y$.
\end{definition}

Una variedad admite infinitas conexiones, pero una m\'etrica distingue una
can\'onica.

\begin{theorem}[Teorema fundamental de la geometr\'ia riemanniana]
\label{thm:geo:levi-civita}
En una variedad riemanniana $(\M, g)$ existe una \'unica conexi\'on af\'in
$\nabla$, la \emph{conexi\'on de Levi--Civita}, que es \emph{compatible con la
m\'etrica}, $X g(Y,Z) = g(\nabla_X Y, Z) + g(Y, \nabla_X Z)$, y \emph{libre de
torsi\'on}, $\nabla_X Y - \nabla_Y X = [X,Y]$. En coordenadas locales queda
determinada por los \emph{s\'imbolos de Christoffel}
\[
  \Gamma^k_{ij} = \tfrac{1}{2} g^{kl}\!\left(
    \partial_j g_{il} + \partial_i g_{jl} - \partial_l g_{ij}\right).
\]
\end{theorem}

La unicidad se sigue de la \emph{f\'ormula de Koszul}, obtenida al expandir
c\'iclicamente la compatibilidad m\'etrica y usar la ausencia de torsi\'on para
cancelar los t\'erminos cruzados; como $g$ es no degenerada, la f\'ormula
determina $\nabla_X Y$ de forma \'unica, y una verificaci\'on directa muestra que
satisface los axiomas. Geom\'etricamente, las dos propiedades que la definen
afirman que el transporte paralelo preserva longitudes y \'angulos, y que los
paralelogramos infinitesimales cierran --- el transporte no introduce rotaciones
espurias.

La conexi\'on de Levi--Civita permite comparar las segundas derivadas
covariantes $\nabla_X \nabla_Y Z$ y $\nabla_Y \nabla_X Z$. En el espacio
euclidiano coinciden; en una variedad curva difieren, y la discrepancia es la
curvatura.

\begin{definition}[Tensor de curvatura de Riemann]
\label{def:geo:riemann}
El \emph{tensor de curvatura de Riemann} de $(\M, g)$ es
\[
  R(X,Y)Z = \nabla_X \nabla_Y Z - \nabla_Y \nabla_X Z - \nabla_{[X,Y]} Z .
\]
Es $C^\infty(\M)$-multilineal, de modo que su valor en $p$ depende solo de
$X, Y, Z$ en $p$. Una variedad riemanniana es localmente isom\'etrica a $\R^n$
si y solo si $R \equiv 0$.
\end{definition}

\begin{definition}[Curvatura seccional, de Ricci y escalar]
\label{def:geo:sectional}
Para un $2$-plano $\sigma \subset \TpM$ generado por $u, v$ independientes, la
\emph{curvatura seccional} es
\[
  K(\sigma) = \frac{\inner{R(u,v)v}{u}_g}{\norm{u}^2\norm{v}^2 - \inner{u}{v}_g^2},
\]
la curvatura de la superficie generada por geod\'esicas en direcciones de
$\sigma$. Contrayendo $R$ se obtiene el \emph{tensor de Ricci}
$R_{ij} = \sum_k \inner{R(\partial_i,\partial_k)\partial_j}{\partial_k}_g$, y
trazando de nuevo la \emph{curvatura escalar} $\mathcal{R} = g^{ij}R_{ij}$.
\end{definition}

La curvatura seccional controla c\'omo se desv\'ia la geometr\'ia de los
postulados de Euclides (\Cref{fig:geo:curvature}). En un tri\'angulo geod\'esico
de \'area $A$, la suma de \'angulos es $\alpha + \beta + \gamma = \pi + K\,A$; un
disco geod\'esico de radio $r$ tiene \'area que crece m\'as r\'apido que
$\pi r^2$ cuando $K < 0$ y m\'as lento cuando $K > 0$; y la separaci\'on
$\xi(t)$ entre geod\'esicas inicialmente paralelas obedece la ecuaci\'on de
Jacobi $\ddot\xi + K\xi = 0$, cuyas soluciones son oscilatorias ($K > 0$),
constantes ($K = 0$) o exponenciales ($K < 0$). Este \'ultimo hecho explica por
qu\'e el espacio hiperb\'olico, con su crecimiento exponencial de volumen,
embebe \'arboles con baja distorsi\'on --- propiedad que explotamos en la
\Cref{subsec:geo:specific}.

\begin{figure}[h]
\centering
\begin{tikzpicture}[scale=1.0]
% --- Curvatura negativa (silla) ---
\begin{scope}[shift={(-5.2,0)}]
  \shade[ball color=gdlred!18, opacity=0.85]
    (-1.6,0.6) .. controls (-0.6,1.5) and (0.6,-0.5) .. (1.6,0.6)
    .. controls (1.0,-0.2) and (0.4,-1.2) .. (0,-1.0)
    .. controls (-0.4,-1.2) and (-1.0,-0.2) .. (-1.6,0.6);
  \draw[gdlred!50, thin] (-1.4,0.5) .. controls (0,1.0) .. (1.4,0.5);
  \draw[gdlred!50, thin] (-0.9,-0.7) .. controls (0,-0.3) .. (0.9,-0.7);
  \node[below, font=\small\bfseries, text=gdlred!70!black] at (0,-1.6) {$K<0$};
  \node[below, font=\footnotesize, text=gdlred!70!black] at (0,-2.0) {divergen};
\end{scope}
% --- Curvatura nula (plano) ---
\begin{scope}[shift={(0,0)}]
  \fill[gdlgreen!15] (-1.6,-1.0) -- (1.6,-1.0) -- (1.6,1.0) -- (-1.6,1.0) -- cycle;
  \foreach \x in {-1.2,-0.6,0,0.6,1.2}
    \draw[gdlgreen!45, thin] (\x,-1.0) -- (\x,1.0);
  \foreach \y in {-0.6,0,0.6}
    \draw[gdlgreen!45, thin] (-1.6,\y) -- (1.6,\y);
  \node[below, font=\small\bfseries, text=gdlgreen!60!black] at (0,-1.6) {$K=0$};
  \node[below, font=\footnotesize, text=gdlgreen!60!black] at (0,-2.0) {paralelas};
\end{scope}
% --- Curvatura positiva (esfera) ---
\begin{scope}[shift={(5.2,0)}]
  \shade[ball color=gdlblue!22, opacity=0.9] (0,0) circle (1.3cm);
  \draw[gdlblue!50, thin] (0,0) circle (1.3cm);
  \draw[gdlblue!50, thin] (-1.3,0) .. controls (0,0.35) .. (1.3,0);
  \draw[gdlblue!50, thin] (0,1.3) .. controls (0.35,0) .. (0,-1.3);
  \node[below, font=\small\bfseries, text=gdlblue!70!black] at (0,-1.6) {$K>0$};
  \node[below, font=\footnotesize, text=gdlblue!70!black] at (0,-2.0) {convergen};
\end{scope}
\end{tikzpicture}
\caption[Comparaci\'on de curvaturas]{Los tres reg\'imenes de curvatura
seccional. Las geod\'esicas que parten paralelas divergen en una superficie de
curvatura negativa ($K<0$, izquierda), permanecen paralelas en un plano
($K=0$, centro) y convergen en una de curvatura positiva ($K>0$, derecha).}
\label{fig:geo:curvature}
\end{figure}

\section{Geod\'esicas y mapas exponencial y logar\'itmico}
\label{sec:geo:geodesics}

Las \emph{geod\'esicas} son las curvas que minimizan localmente la longitud,
generalizando las rectas al espacio curvo. La conexi\'on de Levi--Civita aporta
la derivada covariante que las formaliza: una geod\'esica es una curva cuyo
vector tangente es paralelo a lo largo de s\'i misma, de modo que ``no acelera''
respecto de la geometr\'ia.

\begin{definition}[Geod\'esica]
\label{def:geo:geodesic}
Una curva $\gamma \colon I \to \M$ es una \emph{geod\'esica} si
$\nabla_{\dot\gamma}\dot\gamma = 0$, lo que en coordenadas locales es el sistema
de segundo orden
\[
  \ddot\gamma^k + \Gamma^k_{ij}\, \dot\gamma^i \dot\gamma^j = 0 ,
\]
con la convenci\'on de suma de Einstein.
\end{definition}

Las geod\'esicas admiten tres lecturas complementarias: como extremales del
funcional de longitud (principio variacional), como trayectorias de una
part\'icula libre (inercia) y como caminos localmente m\'as cortos. Son rectas
en $\R^n$, c\'irculos m\'aximos en $S^2$ y arcos de circunferencia que cortan el
borde ortogonalmente en el modelo de Poincar\'e de $\mathbb{H}^n$.

\begin{definition}[Mapas exponencial y logar\'itmico]
\label{def:geo:exp-log}
El \emph{mapa exponencial} $\exp_p \colon \TpM \to \M$ env\'ia $v$ a
$\gamma_v(1)$, donde $\gamma_v$ es la \'unica geod\'esica con $\gamma_v(0) = p$ y
$\dot\gamma_v(0) = v$ --- el punto alcanzado al recorrer la geod\'esica en
direcci\'on $v$ durante tiempo unitario. Su inversa local, el \emph{mapa
logar\'itmico} $\log_p \colon U \to \TpM$, se define en un entorno $U$ de $p$
donde $\exp_p$ es un difeomorfismo: $\log_p(q)$ es el ``vector que apunta de $p$
a $q$''.
\end{definition}

\begin{proposition}[Propiedades del mapa exponencial]
\label{prop:geo:exp}
Para $p \in \M$: (i) $\exp_p(0) = p$; (ii) $\exp_p(tv) = \gamma_v(t)$; (iii)
$d(\exp_p)_0 = \operatorname{id}_{\TpM}$; (iv) $\exp_p$ es un difeomorfismo de un
entorno de $0 \in \TpM$ en un entorno de $p$; y (v)
$d_g(p, \exp_p(v)) = \norm{v}_g$ para $v$ suficientemente peque\~no.
\end{proposition}

La propiedad (iii) se sigue de derivar $t \mapsto \exp_p(tv)$ en $t = 0$; el
teorema de la funci\'on inversa da entonces (iv). La propiedad (v) se cumple
porque la conexi\'on de Levi--Civita mantiene constante $\norm{\dot\gamma_v}_g$,
de modo que la geod\'esica $\gamma_v|_{[0,1]}$ tiene longitud $\norm{v}_g$ y es
minimizante para $v$ peque\~no. La identidad $d_g(p,q) = \norm{\log_p(q)}_g$ es
exactamente el puente entre la distancia intr\'inseca y el espacio tangente
(\Cref{fig:geo:exp-log}).

\begin{figure}[h]
\centering
\begin{tikzpicture}[>=Stealth, scale=1.0]
% Variedad (esfera)
\shade[ball color=gdlblue!18, opacity=0.9] (0,0) circle (2.6cm);
\draw[thick, gdlblue!45] (0,0) circle (2.6cm);
\node[gdlblue!60!black, font=\small] at (-2.1,-2.0) {$\mathcal{M}$};
% Punto p
\coordinate (p) at (0,2.6);
\fill[gdlred!75] (p) circle (2.5pt);
\node[above right, font=\small] at (p) {$p$};
% Plano tangente
\begin{scope}[shift={(0,2.6)}]
  \fill[gdlgreen!16] (-2.2,-0.25) -- (2.2,-0.25) -- (2.2,1.4) -- (-2.2,1.4) -- cycle;
  \draw[gdlgreen!50] (-2.2,0) -- (2.2,0);
  \node[gdlgreen!55!black, font=\small] at (2.4,1.0) {$T_p\mathcal{M}$};
  \draw[->, very thick, gdlorange!80] (0,0) -- (1.3,0.7) node[above, font=\small] {$v$};
\end{scope}
% Geod\'esica
\draw[very thick, gdlred!70, dashed] (0,2.6) arc (90:33:2.6cm);
% Punto q
\coordinate (q) at ({2.6*cos(33)},{2.6*sin(33)});
\fill[gdlpurple!75] (q) circle (2.5pt);
\node[right, font=\small, gdlpurple!75] at (q) {$q=\exp_p(v)$};
% Mapas
\draw[->, very thick, gdlorange!80] (1.3,3.3) .. controls (2.3,2.8) and (2.6,2.0) .. (q);
\node[gdlorange!80, font=\small] at (3.0,3.0) {$\exp_p$};
\draw[->, very thick, gdlpurple!75, dashed] ($(q)+(-0.25,0.3)$) .. controls (2.2,2.5) and (1.7,3.4) .. (1.0,3.5);
\node[gdlpurple!75, font=\small] at (2.6,3.7) {$\log_p$};
% Nota
\node[font=\footnotesize, text width=4.4cm, align=left] at (-4.6,0.2) {
  $\exp_p(v)=\gamma_v(1)$, $\gamma_v$ geod\'esica con $\dot\gamma_v(0)=v$\\[2pt]
  $\log_p=\exp_p^{-1}$ (local)\\[2pt]
  $d_g(p,q)=\norm{\log_p(q)}_g$};
\end{tikzpicture}
\caption[Mapas exponencial y logar\'itmico]{El mapa exponencial env\'ia un
vector tangente $v \in T_p\mathcal{M}$ al extremo de la geod\'esica que genera;
el mapa logar\'itmico es su inversa local. La distancia geod\'esica es igual a la
norma en el espacio tangente de $\log_p(q)$.}
\label{fig:geo:exp-log}
\end{figure}

\paragraph{Coordenadas normales.}
Elegir una base ortonormal $\{e_1, \dots, e_m\}$ de $\TpM$ y definir
$\phi(\exp_p(v^i e_i)) = (v^1, \dots, v^m)$ da las \emph{coordenadas normales
geod\'esicas} centradas en $p$. En ellas, las geod\'esicas por $p$ son rectas
por el origen, $g_{ij}(p) = \delta_{ij}$, los s\'imbolos de Christoffel se anulan
en $p$ y tambi\'en las primeras derivadas de la m\'etrica. El t\'ermino de
segundo orden revela la curvatura:
\[
  g_{ij}(x) = \delta_{ij} - \tfrac{1}{3} R_{ikjl}(p)\, x^k x^l + O(\norm{x}^3) .
\]
As\'i, las coordenadas normales \emph{linealizan} la variedad: dentro de una
bola geod\'esica de radio $r$ parece euclidiana salvo correcciones $O(r^2)$
gobernadas por la curvatura. Esta es precisamente la carta en la que operan las
convoluciones geod\'esicas.

\begin{remark}[La ambig\"uedad gauge]
\label{rem:geo:gauge}
La base $\{e_i\}$ no es can\'onica: cualquier otra base ortonormal difiere por
una rotaci\'on $O \in O(m)$. Esta \emph{libertad gauge} rota la carta de
coordenadas normales, de modo que un kernel de convoluci\'on expresado en ellas
se transforma como $k'(r,\theta) = k(r, \theta - \alpha)$. Hacer que una red sea
independiente de la elecci\'on exige restricciones de equivarianza sobre el
kernel --- el tema del \Cref{ch:gauges}.
\end{remark}

\section{Convoluci\'on geod\'esica}
\label{sec:geo:convolution}

Las convoluciones de grupo (\Cref{ch:groups}) explotan una simetr\'ia global,
pero muchos dominios del mundo real --- superficies 3D, l\'aminas corticales,
espacios de configuraci\'on --- carecen de tal simetr\'ia: su geometr\'ia local
var\'ia punto a punto, no hay una traslaci\'on globalmente consistente ni un
sistema de coordenadas global can\'onico. El mapa exponencial elimina estos
obst\'aculos localmente, y las \emph{redes neuronales convolucionales
geod\'esicas}~\citep{masci2015geodesic} fueron la primera extensi\'on
sistem\'atica de la convoluci\'on a variedades riemannianas generales.

\subsection{Convoluci\'on en cartas geod\'esicas}
\label{subsec:geo:gcnn}

La idea es construir, alrededor de cada punto $p$, una carta local mediante el
mapa exponencial y realizar una convoluci\'on en coordenadas polares dentro de
ella. Una se\~nal $f \colon \M \to \R$ se convoluciona con un kernel aprendible
$k$ muestreando la variedad a lo largo de rayos geod\'esicos:
\begin{equation}
\label{eq:geo:gcnn}
  (f *_g k)(p) = \sum_{i,j} k(r_i, \theta_j)\;
    f\!\left(\exp_p(r_i\, e_{\theta_j})\right)\; J_p(r_i, \theta_j),
\end{equation}
donde $e_{\theta_j} = \cos\theta_j\, e_1 + \sin\theta_j\, e_2$ es la direcci\'on
angular, $(r_i, \theta_j)$ discretiza un disco geod\'esico en coordenadas
polares y $J_p(r, \theta) = \sqrt{\det g_{\exp_p(r e_\theta)}}$ es un factor
jacobiano que corrige la distorsi\'on entre el parche tangente y la variedad.
Por la expansi\'on en coordenadas normales del entorno de \Cref{rem:geo:gauge},
\[
  J_p(r, \theta) = 1 - \tfrac{1}{6} K(p)\, r^2 + O(r^3),
\]
de modo que la correcci\'on es despreciable para parches peque\~nos y crece con
la curvatura de Gauss $K(p)$. En la pr\'actica, construir cada carta requiere
elegir una base tangente ortonormal, mapear una rejilla polar a trav\'es de
$\exp_p$ e interpolar los valores de las caracter\'isticas en los puntos
resultantes (\Cref{fig:geo:convolution}).

\begin{figure}[h]
\centering
\begin{tikzpicture}[>=Stealth, scale=1.0]
% Parche de superficie curva
\shade[ball color=gdlorange!20, opacity=0.9]
  (-2.6,0) .. controls (-2.0,1.4) and (2.0,1.4) .. (2.6,0)
  .. controls (2.0,-1.2) and (-2.0,-1.2) .. (-2.6,0);
\draw[gdlorange!55, thick]
  (-2.6,0) .. controls (-2.0,1.4) and (2.0,1.4) .. (2.6,0)
  .. controls (2.0,-1.2) and (-2.0,-1.2) .. (-2.6,0);
\node[gdlorange!60!black, font=\small] at (-2.4,-1.1) {$\mathcal{M}$};
% Punto central p
\coordinate (p) at (0,0.25);
% Rejilla polar geod\'esica (anillos + rayos)
\foreach \rr in {0.45,0.9,1.35}
  \draw[gdlblue!55, thin] (p) ellipse ({\rr} and {0.55*\rr});
\foreach \aa in {0,45,90,135,180,225,270,315}
  \draw[gdlblue!40, thin] (p) -- ($(p)+({1.35*cos(\aa)}, {0.55*1.35*sin(\aa)})$);
\fill[gdlred!80] (p) circle (2.5pt);
\node[above, font=\small] at (p) {$p$};
% Un vecino
\coordinate (q) at ($(p)+({1.35*cos(35)}, {0.55*1.35*sin(35)})$);
\fill[gdlgreen!70] (q) circle (2pt);
\node[right, font=\small, gdlgreen!60!black] at (q) {$q_j=\exp_p(r_i e_{\theta_j})$};
% etiqueta
\node[font=\footnotesize, text=gdlblue!60!black] at (0,-1.5) {bola geod\'esica $B_r(p)$};
\end{tikzpicture}
\caption[Convoluci\'on geod\'esica]{Convoluci\'on geod\'esica. Una rejilla polar
de radios $r_i$ y \'angulos $\theta_j$ se dispone en el espacio tangente en $p$
y se mapea sobre la bola geod\'esica $B_r(p)$ a trav\'es de $\exp_p$; los pesos
del kernel $k(r_i,\theta_j)$ se aplican a las caracter\'isticas interpoladas,
corregidas por el jacobiano $J_p$.}
\label{fig:geo:convolution}
\end{figure}

Las convoluciones geod\'esicas fueron un avance conceptual, pero tienen
limitaciones claras. La construcci\'on depende de una elecci\'on de base
tangente en cada punto, que puede ser inconsistente entre vecinos (la
ambig\"uedad gauge de \Cref{rem:geo:gauge}); calcular cartas geod\'esicas y
jacobianos es costoso en mallas de alta resoluci\'on; y el mapa exponencial
requiere suficiente suavidad, por lo que los dominios singulares son
problem\'aticos. Los \emph{descriptores de difusi\'on
anisotr\'opica}~\citep{Boscaini2016AnisotropicDD} abordan la carencia an\'aloga
de los descriptores espectrales isotr\'opicos al reemplazar el n\'ucleo del calor
por una difusi\'on dependiente de la direcci\'on
$\partial_t u = -\divergence(C \nabla u)$ cuyo tensor de conductividad $C$ se
adapta a la geometr\'ia local, capturando la orientaci\'on de caracter\'isticas
alargadas.

\subsection{Un marco general: MoNet}
\label{subsec:geo:monet}

La dependencia de un sistema de coordenadas expl\'icito motiva una formulaci\'on
m\'as abstracta. MoNet~\citep{DBLP:journals/corr/MontiBMRSB16} reemplaza las
cartas geod\'esicas por \emph{pseudo-coordenadas}: una aplicaci\'on
$u \colon \M \times \M \to \R^d$ que codifica la relaci\'on geom\'etrica local
entre un punto $x$ y sus vecinos, es local (se anula fuera de un entorno),
invariante bajo isometr\'ias y suave. La convoluci\'on se vuelve una agregaci\'on
ponderada sobre un entorno $\mathcal{N}(x)$,
\begin{equation}
\label{eq:geo:monet}
  (f * w)(x) = \int_{\mathcal{N}(x)} w_\theta\big(u(x,y)\big)\, f(y)\, dV_g(y),
\end{equation}
donde el kernel $w_\theta$ es una mezcla aprendible de funciones base
param\'etricas, t\'ipicamente gaussianas,
\[
  w_\theta(u) = \sum_{j=1}^J \alpha_j\,
    \exp\!\left(-\tfrac{1}{2}(u - \mu_j)^\top \Sigma_j^{-1} (u - \mu_j)\right),
\]
con pesos $\alpha_j$, medias $\mu_j$ y covarianzas $\Sigma_j$ aprendibles. En una
malla discreta la integral se vuelve
$(f * w)(x_i) = \sum_{j \in \mathcal{N}(i)} w_\theta(u(x_i,x_j))\, f(x_j)\, A_j$
con pesos de \'area local $A_j$.

\begin{theorem}[Aproximaci\'on universal de los kernels de MoNet]
\label{thm:geo:monet}
Para toda funci\'on continua de soporte compacto $k \colon \R^d \to \R$ y todo
$\varepsilon > 0$ existen $J$ y par\'ametros
$\{\alpha_j, \mu_j, \Sigma_j\}_{j=1}^J$ con
$\sup_{u}\, \abs{k(u) - w_\theta(u)} < \varepsilon$.
\end{theorem}

Esto se sigue de la densidad de las mezclas gaussianas junto con la compacidad
del soporte, y convierte a MoNet en una generalizaci\'on genuina: distintas
elecciones de $u$ recuperan modelos anteriores. Tomando
$u(x,y) = (d_g(x,y), \angle(x,y))$ (distancia geod\'esica y \'angulo polar) se
reproducen las convoluciones geod\'esicas~\eqref{eq:geo:gcnn}; tomando
coordenadas basadas en el grado $u(x,y) = (1/\sqrt{\deg x},\, 1/\sqrt{\deg y})$
se recuperan las redes convolucionales de grafos (\Cref{ch:graphs}); y
coordenadas de B-splines recuperan las convoluciones basadas en splines. As\'i,
MoNet unifica las convoluciones en rejillas, grafos y variedades bajo una \'unica
receta de kernel aprendible, y es a su vez una instancia de la plantilla de paso
de mensajes del \Cref{ch:graphs}.

\subsection{Transporte paralelo y variedades espec\'ificas}
\label{subsec:geo:specific}

Comparar vectores de caracter\'isticas que viven en espacios tangentes distintos
exige llevar uno al otro. La conexi\'on de Levi--Civita proporciona el
\emph{transporte paralelo} $\Pi_{p \to q}^\gamma \colon \TpM \to T_q\M$ a lo
largo de la geod\'esica $\gamma$ que une $p$ y $q$; es una isometr\'ia lineal,
preserva normas y \'angulos, y se compone a lo largo de caminos concatenados. La
holonom\'ia alrededor de lazos cerrados est\'a controlada por la curvatura (el
teorema de Ambrose--Singer identifica el \'algebra de Lie de holonom\'ia con los
operadores de curvatura), de modo que el transporte paralelo codifica la
geometr\'ia intr\'inseca completa. Las \emph{convoluciones por transporte
paralelo}~\citep{schonsheck2018parallel} lo explotan para definir una
convoluci\'on libre de coordenadas que transporta el kernel de forma consistente
por la superficie, evitando la ambig\"uedad gauge de las cartas geod\'esicas.
Estas ideas conectan directamente con las redes equivariantes gauge del
\Cref{ch:gauges}.

Los mapas exponencial y logar\'itmico tambi\'en permiten reformular bloques
cl\'asicos de manera intr\'inseca. Una \emph{conexi\'on residual riemanniana}
reemplaza la actualizaci\'on euclidiana $x_{\ell+1} = x_\ell + F_\ell(x_\ell)$
por
\[
  x_{\ell+1} = \exp_{x_\ell}\!\big(F_\ell(x_\ell)\big),
    \qquad F_\ell(x_\ell) \in T_{x_\ell}\M,
\]
que mantiene el estado sobre la variedad, y la interpolaci\'on geod\'esica
$\gamma(t) = \exp_p(t \log_p(q))$ generaliza la interpolaci\'on lineal.

Cuando los datos tienen una geometr\'ia conocida puede especializarse la
variedad. Los datos jer\'arquicos --- \'arboles, taxonom\'ias, ontolog\'ias ---
se embeben de forma natural en el \emph{espacio hiperb\'olico}: como el volumen
de una bola hiperb\'olica crece exponencialmente con su radio (el r\'egimen
$K < 0$ de \Cref{fig:geo:curvature}), un \'arbol cuyo n\'umero de nodos crece
exponencialmente con la profundidad encaja con distorsi\'on acotada, algo
imposible en cualquier espacio euclidiano de dimensi\'on fija. En la bola de
Poincar\'e se reemplaza la suma de vectores por la operaci\'on de M\"obius y la
distancia euclidiana por la distancia hiperb\'olica; las \emph{redes de
atenci\'on hiperb\'olica}~\citep{gulcehre2019hyperbolic} ponderan la atenci\'on
por $\exp(-d_{\mathbb{H}}(x,y)/\tau)$ con la distancia hiperb\'olica
$d_{\mathbb{H}}$. El principio unificador es el mismo en todos los casos:
respetar la geometr\'ia del espacio de datos al dise\~nar cada componente, de
modo que la equivarianza bajo el grupo de isometr\'ias haga intr\'insecas las
representaciones aprendidas.

\subsection{Casos de estudio}
\label{subsec:geo:cases}

Tres aplicaciones muestran a los m\'etodos geod\'esicos superando a las
referencias euclidianas.

\paragraph{An\'alisis de formas 3D (FAUST).}
El benchmark FAUST contiene mallas trianguladas de cuerpos humanos (unos $6890$
v\'ertices cada una) en poses variadas, con fuertes variaciones de curvatura.
Como las distancias geod\'esicas son invariantes bajo las deformaciones (casi
isom\'etricas) de un cuerpo en movimiento, las convoluciones geod\'esicas y MoNet
reconocen correspondencias y poses con mucha m\'as precisi\'on que una red
euclidiana aplicada a una proyecci\'on plana, que es sensible al embebido. Las
regiones de alta curvatura, como las articulaciones, son las que m\'as se
benefician, y la correcci\'on jacobiana es cr\'itica cerca de ellas.

\paragraph{Im\'agenes omnidireccionales ($360^\circ$).}
Una imagen esf\'erica proyectada a un plano (equirectangular) sufre una severa
distorsi\'on polar --- las \'areas cercanas a los polos se estiran en un orden de
magnitud, rompiendo la invarianza a traslaci\'on que asume una convoluci\'on
est\'andar. Tratar la imagen como una se\~nal en $S^2$ y convolucionar de forma
geod\'esica (por ejemplo, sobre una discretizaci\'on icosa\'edrica) mantiene la
distorsi\'on casi uniforme y preserva las formas locales, mejorando la
clasificaci\'on de escenas y la segmentaci\'on sem\'antica.

\paragraph{Superficies m\'edicas (corteza cerebral).}
La superficie cortical es una variedad de g\'enero 0 con surcos profundos y
variaci\'on extrema de curvatura, y presenta gran variabilidad anat\'omica entre
sujetos. Convolucionar de forma intr\'inseca sobre la superficie cortical, en
lugar de sobre una representaci\'on volum\'etrica o proyectada, respeta el
patr\'on de plegamiento y produce caracter\'isticas anat\'omicamente
consistentes para tareas como la parcelaci\'on en grandes cohortes.

\subsection{Geometr\'ia hiperb\'olica y atenci\'on}
\label{subsec:geo:hyperbolic}

El espacio hiperb\'olico merece una mirada m\'as detenida, porque es la variedad
que m\'as a menudo se especializa en la pr\'actica y porque conecta el
aprendizaje geod\'esico con los mecanismos de atenci\'on del \Cref{ch:graphs}.
Su rasgo definitorio, el r\'egimen $K = -1$ de \Cref{fig:geo:curvature}, es el
crecimiento exponencial del volumen: una bola geod\'esica de radio $r$ en
$\mathbb{H}^n$ tiene volumen de orden $e^{(n-1)r}$, que iguala la explosi\'on
combinatoria de los \'arboles, cuyo n\'umero de nodos crece exponencialmente con
la profundidad. Se usan dos modelos equivalentes. En el \emph{modelo del
hiperboloide (de Lorentz)} $\mathbb{H}^n = \{x \in \R^{n+1} :
\inner{x}{x}_{\mathcal L} = -1,\, x_0 > 0\}$, donde
$\inner{x}{y}_{\mathcal L} = -x_0 y_0 + \sum_{i\ge 1} x_i y_i$ es la forma de
Minkowski; en la \emph{bola de Poincar\'e} $\mathbb{D}^n$ de
\Cref{ex:geo:metrics}, conforme a la m\'etrica euclidiana. Ambos se relacionan
por un difeomorfismo isom\'etrico, de modo que se elige el m\'as conveniente para
el c\'omputo.

\begin{definition}[Distancia hiperb\'olica]
\label{def:geo:hyperbolic-distance}
En el modelo del hiperboloide, la distancia geod\'esica entre
$x, y \in \mathbb{H}^n$ es
\[
  d_{\mathbb{H}}(x,y) = \operatorname{arccosh}\big(-\inner{x}{y}_{\mathcal L}\big),
\]
y en la bola de Poincar\'e la expresi\'on equivalente es
\[
  d_{\mathbb{D}}(u,v) = \operatorname{arccosh}\!\left(
    1 + 2\,\frac{\norm{u-v}^2}{(1 - \norm{u}^2)(1 - \norm{v}^2)}\right).
\]
Es una m\'etrica, es sim\'etrica, satisface la desigualdad triangular y es
invariante bajo el grupo de isometr\'ias $\mathrm{SO}^+(1,n)$.
\end{definition}

La bola de Poincar\'e admite tambi\'en an\'alogos algebraicos de las operaciones
lineales, las \emph{operaciones de M\"obius}: la suma de M\"obius
$u \oplus v$ que generaliza la suma de vectores, el producto escalar de M\"obius
$r \otimes u = \tanh\!\big(r\,\tanh^{-1}\norm{u}\big)\,u/\norm{u}$ y el mapa
exponencial hiperb\'olico
$\exp^{\mathbb{H}}_x(v) = x \oplus \tanh\!\big(\tfrac{1}{2}\lambda_x\norm{v}\big)
\,v/\norm{v}$ con factor conforme $\lambda_x = 2/(1-\norm{x}^2)$. Estas son las
instancias en la bola de Poincar\'e de la operaci\'on de M\"obius y del mapa
exponencial referidos en \Cref{subsec:geo:specific}.

\begin{theorem}[Embebido de \'arboles con baja distorsi\'on]
\label{thm:geo:tree-embedding}
Todo \'arbol finito $T$ admite un embebido en el plano hiperb\'olico
$\mathbb{H}^2$ con distorsi\'on arbitrariamente peque\~na: para todo
$\varepsilon > 0$ existe una aplicaci\'on $f \colon V(T) \to \mathbb{H}^2$ con
\[
  d_T(u,v) \le d_{\mathbb{H}}(f(u), f(v)) \le (1+\varepsilon)\, d_T(u,v)
  \qquad \text{para todo } u, v \in V(T),
\]
donde $d_T$ es la distancia en el \'arbol (geod\'esica de grafo)~\citep{sarkar2011low}.
En contraste, embeber un \'arbol binario completo de $n$ v\'ertices en $\R^d$ con
distorsi\'on acotada exige $d = \Omega(\log n)$.
\end{theorem}

La construcci\'on coloca la ra\'iz en el origen y, en cada nodo interno,
distribuye los hijos uniformemente sobre un c\'irculo hiperb\'olico centrado en
el padre; como el per\'imetro de un c\'irculo hiperb\'olico de radio $r$ crece
como $2\pi\sinh r \sim \pi e^r$, siempre hay sitio para separar las ramas sin
distorsi\'on. El crecimiento polin\'omico del volumen $\Theta(r^d)$ de $\R^d$ lo
impide en dimensi\'on baja, que es justo la raz\'on por la que los datos
jer\'arquicos --- WordNet, taxonom\'ias, grafos de conocimiento --- se embeben en
la bola de Poincar\'e, con la norma $\norm{f(e)}$ codificando la profundidad de
una entidad en la jerarqu\'ia (\Cref{fig:geo:hyperbolic}).

\begin{figure}[h]
\centering
\begin{tikzpicture}[>=Stealth, scale=1.0]
% Borde del disco de Poincar\'e
\draw[thick, gdlpurple!55] (0,0) circle (2.6cm);
\fill[gdlpurple!6] (0,0) circle (2.6cm);
\node[gdlpurple!60!black, font=\small] at (-2.0,-2.3) {$\mathbb{D}^2$};
% ra\'iz
\coordinate (r) at (0,0);
\fill[gdlred!80] (r) circle (2.5pt);
\node[above, font=\footnotesize] at (r) {ra\'iz};
% nivel 1
\foreach \a in {90,210,330}{
  \coordinate (c\a) at ({1.2*cos(\a)},{1.2*sin(\a)});
  \draw[gdlblue!55, thick] (r) -- (c\a);
  \fill[gdlblue!70] (c\a) circle (2pt);
}
% nivel 2 alrededor de cada nodo de nivel 1 (empujado hacia el borde)
\foreach \a in {90,210,330}{
  \foreach \d in {-28,28}{
    \coordinate (l\a\d) at ({2.25*cos(\a+\d)},{2.25*sin(\a+\d)});
    \draw[gdlgreen!55] (c\a) -- (l\a\d);
    \fill[gdlgreen!70] (l\a\d) circle (1.6pt);
  }
}
\node[font=\footnotesize, text width=4.0cm, align=left] at (4.7,0) {
  crecimiento exponencial\\ del volumen $\Rightarrow$ las ramas\\ se separan sin\\
  distorsi\'on (\Cref{thm:geo:tree-embedding})};
\end{tikzpicture}
\caption[Embebido hiperb\'olico de \'arboles]{Una jerarqu\'ia embebida en el
disco de Poincar\'e $\mathbb{D}^2$. La ra\'iz se sit\'ua en el centro y los
niveles sucesivos se empujan hacia el borde; el crecimiento exponencial del
volumen hiperb\'olico da espacio para separar las ramas con distorsi\'on
evanescente, algo imposible en un espacio euclidiano de dimensi\'on fija.}
\label{fig:geo:hyperbolic}
\end{figure}

Cuando los datos son hiperb\'olicos, el propio mecanismo de atenci\'on puede
trasladarse a $\mathbb{H}^n$. Las \emph{redes de atenci\'on
hiperb\'olica}~\citep{gulcehre2019hyperbolic} calculan los pesos de atenci\'on a
partir de la distancia hiperb\'olica,
\[
  \alpha_{ij} \;\propto\; \exp\!\big({-}\beta\, d_{\mathbb{H}}(q_i, k_j)\big),
\]
y agregan los valores con una media hiperb\'olica --- el punto medio de Einstein
en el modelo de Klein o la \emph{media de Fr\'echet}
$\operatorname*{arg\,min}_{x \in \mathbb{H}^d} \sum_i \alpha_{ij}\,
d_{\mathbb{H}}^2(x, v_j)$, calculada con unos pocos pasos de gradiente
riemanniano usando los mapas hiperb\'olicos $\exp$ y $\log$. Como
$d_{\mathbb{H}}$ es invariante bajo $\mathrm{SO}^+(1,n)$
(\Cref{def:geo:hyperbolic-distance}), los pesos de atenci\'on son invariantes y
la agregaci\'on por media de Fr\'echet es equivariante bajo las isometr\'ias
hiperb\'olicas, de modo que toda la capa respeta el grupo de isometr\'ias del
espacio de datos --- el mismo principio de equivarianza del \Cref{ch:priors} que
gobierna los filtros espectrales de m\'as abajo. Una construcci\'on sim\'etrica
en la esfera $S^2$, usando la distancia esf\'erica
$d_{S^2}(x,y) = \arccos\inner{x}{y}$ en lugar de $d_{\mathbb{H}}$, da la
atenci\'on esf\'erica para datos direccionales y omnidireccionales, y el
transporte paralelo (\Cref{subsec:geo:specific}) proporciona la atenci\'on por
transporte que lleva consultas y claves a un espacio tangente com\'un antes de
compararlas.

\section{Geod\'esicas en grafos}
\label{sec:geo:graphs}

Los caminos m\'as cortos en grafos son la discretizaci\'on natural de las
geod\'esicas continuas, y precisan el v\'inculo, anticipado en \Cref{ch:graphs},
entre la distancia combinatoria sobre un grafo de proximidad y la distancia
intr\'inseca sobre la variedad que muestrea.

\begin{definition}[Distancia geod\'esica en grafos]
\label{def:geo:graph-distance}
Para un grafo $G = (V, E)$, la \emph{distancia geod\'esica discreta} entre
v\'ertices $u, v \in V$ es el n\'umero m\'inimo de aristas en un camino que los
une,
\[
  d_G(u,v) = \min\{ k : \exists\, u = v_0, v_1, \dots, v_k = v
    \text{ con } (v_i, v_{i+1}) \in E\}.
\]
En un grafo ponderado se reemplaza el conteo de aristas por la suma de los pesos
a lo largo del camino.
\end{definition}

\begin{definition}[Grafo de vecindad geod\'esico]
\label{def:geo:graph-ball}
Dados $G = (V, E)$ y un radio $r \in \N$, el \emph{$r$-vecindario geod\'esico} de
un v\'ertice $v$ es
\[
  B_r(v) = \{ u \in V : d_G(v, u) \le r\},
\]
el an\'alogo discreto de una bola geod\'esica en una variedad.
\end{definition}

Cuando un grafo $G_n$ se construye muestreando $n$ puntos de una variedad $\M$ y
uniendo puntos cercanos, se cumplen dos convergencias complementarias bajo un
escalamiento apropiado. El laplaciano de grafo normalizado converge al operador
de Laplace--Beltrami (la convergencia \emph{espectral} ya usada en
\Cref{ch:graphs} para justificar el an\'alisis de Fourier intr\'inseco), y la
distancia de camino m\'as corto converge a la distancia geod\'esica (una
convergencia \emph{m\'etrica}, que demostramos ahora). Juntas certifican que un
grafo suficientemente denso aproxima la variedad tanto local como globalmente.

\begin{theorem}[Convergencia de la distancia de camino m\'as corto a la geod\'esica]
\label{thm:geo:graph-convergence}
Sea $(\M, g)$ una variedad riemanniana compacta y conexa de dimensi\'on $d$ con
curvatura seccional acotada, embebida en $\R^N$. Sea
$X_n = \{x_1, \dots, x_n\}$ muestreado i.i.d.\ de una densidad $p$ con
$p(x) \ge p_{\min} > 0$. Constr\'uyase el grafo $G_n = (X_n, E_n)$ con arista
$(x_i,x_j) \in E_n$ si y solo si $\norm{x_i - x_j} \le \varepsilon_n$ y peso
$w_{ij} = \norm{x_i - x_j}$. Si $\varepsilon_n \to 0$ y
$n\varepsilon_n^d / \log n \to \infty$, entonces con alta probabilidad
\[
  \sup_{x_i, x_j \in X_n}
    \big| d_{G_n}(x_i, x_j) - d_g(x_i, x_j)\big| \to 0
  \qquad (n \to \infty).
\]
\end{theorem}

\begin{proof}
Seguimos a~\citet{bernstein2000graph} en cinco pasos.

\textbf{Paso 1: Comparaci\'on m\'etrica local.} Sea $\abs{K} \le \kappa$ la cota
de curvatura seccional. Para $x, y \in \M$ con $d_g(x,y) \le \varepsilon$, la
distancia euclidiana ambiente y la geod\'esica satisfacen
\begin{equation}
\label{eq:geo:local-comparison}
  d_g(x,y)\big(1 - C_\kappa \varepsilon^2\big) \le \norm{x - y} \le d_g(x,y),
\end{equation}
con $C_\kappa > 0$ dependiente solo de $\kappa$ y $d$. La cota superior es
inmediata: el segmento recto en $\R^N$ no es m\'as largo que la geod\'esica. Para
la cota inferior, trabajamos en las coordenadas normales centradas en $x$ de
\Cref{rem:geo:gauge}: por la expansi\'on de la m\'etrica
$g_{ij}(z) = \delta_{ij} - \tfrac13 R_{ikjl}(x)z^k z^l + O(\norm{z}^3)$ y
$\abs{R_{ikjl}} \le C\kappa$, la longitud de la geod\'esica minimizante es
$d_g(x,y) \le \norm{x-y}(1 + C'_\kappa \varepsilon^2)$, que se invierte a la cota
inferior.

\textbf{Paso 2: Recubrimiento probabil\'istico.} Definimos el radio de
recubrimiento $r_n = \max_{q \in \M} \min_{x_i} d_g(q, x_i)$. Afirmamos que, bajo
$n\varepsilon_n^d/\log n \to \infty$, con alta probabilidad
\begin{equation}
\label{eq:geo:covering}
  r_n = O\!\Big(\big(\tfrac{\log n}{n}\big)^{1/d}\Big) = o(\varepsilon_n).
\end{equation}
Como las bolas geod\'esicas tienen volumen
$\operatorname{Vol}(B_g(q,\delta)) \ge c_d\,\delta^d$ para $\delta$ peque\~no, y
$p \ge p_{\min}$, la probabilidad de que ning\'un punto caiga en
$B_g(q,\delta)$ es a lo sumo $\exp(-n\,p_{\min}\,c_d\,\delta^d)$. Una
$\delta$-red de $\M$ tiene $O(\delta^{-d})$ bolas, luego una cota de uni\'on da
$\Pr[r_n > \delta] \le O(\delta^{-d})\exp(-n\,p_{\min}\,c_d\,\delta^d)$.
Eligiendo $\delta = C(\log n / n)^{1/d}$ con $C$ grande, esto tiende a $0$, y la
hip\'otesis de escalamiento da $r_n = o(\varepsilon_n)$.

\textbf{Paso 3: Cota inferior de $d_{G_n}$.} Para un camino \'optimo del grafo
$x_i = z_0, \dots, z_L = x_j$, cada arista cumple
$\norm{z_l - z_{l+1}} \le \varepsilon_n$, luego por el Paso~1
$\norm{z_l - z_{l+1}} \ge d_g(z_l, z_{l+1})(1 - C_\kappa \varepsilon_n^2)$.
Sumando y aplicando la desigualdad triangular de $d_g$
(\Cref{prop:geo:dg-metric}),
\begin{equation}
\label{eq:geo:lower-dGn}
  d_{G_n}(x_i, x_j) \ge (1 - C_\kappa \varepsilon_n^2)\, d_g(x_i, x_j).
\end{equation}

\textbf{Paso 4: Cota superior de $d_{G_n}$.} Sea $\gamma$ la geod\'esica
minimizante de $x_i$ a $x_j$, parametrizada por longitud de arco, y muestreamos
puntos $p_0 = x_i, p_1, \dots$ con paso $h = \varepsilon_n/3$. Por el Paso~2,
$r_n < \varepsilon_n/6$ para $n$ grande, luego cada $p_l$ tiene una muestra $z_l$
con $d_g(z_l, p_l) \le r_n$ (tomando $z_0 = x_i$, \'ultima $= x_j$). Las muestras
consecutivas cumplen
$\norm{z_l - z_{l+1}} \le d_g(z_l, z_{l+1}) \le 2r_n + h < \varepsilon_n$, luego
$(z_l, z_{l+1}) \in E_n$, y
\begin{equation}
\label{eq:geo:upper-dGn}
  d_{G_n}(x_i, x_j) \le \sum_l d_g(z_l, z_{l+1})
    \le \big(\tfrac{L_\gamma}{h} + 1\big)(h + 2r_n)
    = d_g(x_i, x_j) + O\!\Big(d_g(x_i,x_j)\tfrac{r_n}{\varepsilon_n}\Big)
      + O(\varepsilon_n).
\end{equation}

\textbf{Paso 5: Convergencia uniforme.} Combinando
\eqref{eq:geo:lower-dGn}--\eqref{eq:geo:upper-dGn},
\[
  \big| d_{G_n}(x_i,x_j) - d_g(x_i,x_j)\big|
    \le d_g(x_i,x_j)\, O\!\Big(\max\big(\varepsilon_n^2, \tfrac{r_n}{\varepsilon_n}\big)\Big)
      + O(\varepsilon_n).
\]
Como $\M$ es compacta, $d_g \le \operatorname{diam}(\M) < \infty$, y tanto
$\varepsilon_n \to 0$ como $r_n/\varepsilon_n \to 0$ por \eqref{eq:geo:covering},
de modo que el supremo sobre todos los pares tiende a $0$.
\end{proof}

\begin{remark}[Dos convergencias complementarias]
\label{rem:geo:two-convergences}
La convergencia m\'etrica anterior y la convergencia espectral del
\Cref{ch:graphs} son complementarias, no equivalentes. El enunciado espectral
(laplaciano de grafo $\to \Delta_\M$) captura la estructura diferencial local
--- difusi\'on, curvatura, autofunciones --- y requiere el escalamiento m\'as
denso $n\varepsilon^{d+2} \to \infty$ (se necesitan m\'as vecinos para estimar
las segundas derivadas del n\'ucleo del calor); el enunciado m\'etrico
(caminos m\'as cortos $\to d_g$) captura la geometr\'ia global y solo necesita el
escalamiento m\'as d\'ebil de conectividad y cobertura
$n\varepsilon_n^d/\log n \to \infty$. Juntos justifican tratar un grafo
suficientemente denso como un sustituto fiel de la variedad.
\end{remark}

\begin{example}[Aproximaci\'on num\'erica en $S^2$]
\label{ex:geo:sphere-numerical}
En la esfera unitaria la distancia geod\'esica es
$d_g(p,q) = \arccos\inner{p}{q}$ (\Cref{ex:geo:sphere-distance}). Tomemos
$p = (1,0,0)$ y $q = (0,1,0)$, de modo que $d_g(p,q) = \pi/2 \approx 1.571$.
Muestreando $n$ puntos uniformemente en $S^2$ y conectando los que est\'an a
distancia euclidiana menor que $\varepsilon$, la distancia de camino m\'as corto
$d_{G_n}(p,q)$ se aproxima a $\pi/2$ al refinar el muestreo:
\begin{center}
\begin{tabular}{@{}ccc@{}}
\toprule
$n$ (puntos) & $\varepsilon$ & $d_{G_n}(p, q)$ \\
\midrule
$100$ & $0.5$ & $1.62$ \\
$500$ & $0.3$ & $1.59$ \\
$2000$ & $0.2$ & $1.575$ \\
$10000$ & $0.1$ & $1.572$ \\
\bottomrule
\end{tabular}
\end{center}
La convergencia a $d_g = \pi/2$ ilustra \Cref{thm:geo:graph-convergence}: la
distancia discreta sobreestima la geod\'esica en grafos groseros (el camino
zigzaguea entre muestras) y se ajusta al crecer $n$ y reducirse $\varepsilon$ al
ritmo prescrito.
\end{example}

Estas convergencias tienen consecuencias directas para el dise\~no de
arquitecturas. Las codificaciones posicionales construidas a partir de distancias
de camino m\'as corto (como en los \emph{transformers} de grafos) aproximan
distancias geod\'esicas intr\'insecas cuando el grafo muestrea densamente una
variedad; el paso de mensajes basado en difusi\'on captura informaci\'on a escala
$O(\varepsilon)$, mientras que la distancia geod\'esica habilita una atenci\'on de
largo alcance consciente de la geometr\'ia; y el escalamiento
$\varepsilon_n \sim (n/\log n)^{-1/d}$ prescribe con qu\'e rapidez debe reducirse
el radio de conectividad para preservar a la vez la conectividad y la geometr\'ia
local.

\section{M\'etodos espectrales en variedades}
\label{sec:geo:spectral}

La convoluci\'on geod\'esica es una construcci\'on \emph{intr\'inseca pero
espacial}: opera en cartas geod\'esicas. El punto de vista espectral ofrece una
alternativa complementaria, totalmente libre de coordenadas, construida
directamente sobre el operador de Laplace--Beltrami de
\Cref{def:geo:laplace-beltrami}. Es la contraparte en variedades de las
convoluciones espectrales de grafos del \Cref{ch:graphs}, y hace exacta la
analog\'ia entre los mundos continuo y discreto.

\paragraph{La transformada de Fourier en variedades.}
En una variedad compacta, $\Delta_\M$ tiene un espectro discreto
$0 = \lambda_0 \le \lambda_1 \le \cdots$ con autofunciones $\{\varphi_k\}$ que
forman una base ortonormal de $L^2(\M)$ (\Cref{rem:geo:lb-role}). Estas
desempe\~nan el papel de las exponenciales complejas: en $\R^n$ los modos de
Fourier $e^{i\inner{\xi}{x}}$ son exactamente las autofunciones del laplaciano
euclidiano, y en una variedad general las $\varphi_k$ las generalizan. Se define
la \emph{transformada de Fourier en variedades} y su inversa por
\[
  \hat f(\lambda_k) = \inner{f}{\varphi_k} = \int_\M f\, \varphi_k\, dV_g,
  \qquad
  f = \sum_{k} \hat f(\lambda_k)\, \varphi_k .
\]
Los autovalores bajos corresponden a modos suaves y de variaci\'on lenta, y los
altos a modos oscilatorios, igual que las frecuencias bajas y altas en el
an\'alisis de Fourier cl\'asico.

\paragraph{Convoluci\'on espectral e invarianza bajo isometr\'ias.}
Como en una variedad general no hay traslaci\'on, la convoluci\'on se
\emph{define} mediante el teorema de convoluci\'on: un filtro espectral act\'ua
modulando cada coeficiente de Fourier,
\begin{equation}
\label{eq:geo:spectral-conv}
  (g_\eta \star f) = \sum_k \eta(\lambda_k)\, \hat f(\lambda_k)\, \varphi_k,
\end{equation}
donde $\eta$ es una funci\'on de transferencia aprendible del autovalor. Como
$\Delta_\M$ conmuta con las isometr\'ias, tambi\'en lo hace cualquier filtro de
la forma~\eqref{eq:geo:spectral-conv}: la arquitectura es autom\'aticamente
invariante bajo las simetr\'ias de la variedad, realizando el principio de
equivarianza del \Cref{ch:priors} sin restricciones adicionales. Elegir $\eta$
como un polinomio en $\lambda$ hace el filtro \emph{local} (un polinomio de grado
$K$ alcanza solo el vecindario geod\'esico de $K$ saltos) y evita calcular la
descomposici\'on espectral completa --- exactamente el argumento de
localizaci\'on que, en el caso discreto, produjo los filtros de Chebyshev y las
convoluciones de grafos del \Cref{ch:graphs}.

\paragraph{El puente continuo--discreto.}
La construcci\'on espectral es lo que hace de las im\'agenes de variedad y de
grafo dos caras de una misma teor\'ia. La \Cref{tab:geo:analogy} alinea la
correspondencia: el laplaciano de grafo discretiza $\Delta_\M$, sus
autovectores discretizan las autofunciones $\varphi_k$, y las transformadas de
Fourier de grafo y de variedad coinciden en el l\'imite de muestreo. Un filtro
espectral aprendido sobre una malla suficientemente densa aproxima, por tanto, un
operador continuo independiente de la malla --- la propiedad que permite a tales
modelos transferirse entre discretizaciones de la misma forma.

\begin{table}[h]
\centering
\caption{Correspondencia entre el an\'alisis espectral en variedades y en grafos.}
\label{tab:geo:analogy}
\begin{tabular}{@{}p{5.4cm}p{6.0cm}@{}}
\toprule
\textbf{Variedad $(\M, g)$} & \textbf{Grafo $G = (V,E)$} \\
\midrule
Laplace--Beltrami $\Delta_\M$ & Laplaciano de grafo $L = D - W$ \\
\addlinespace
Autofunciones $\varphi_k$ & Autovectores del laplaciano $u_k$ \\
\addlinespace
Espectro $0 = \lambda_0 \le \lambda_1 \le \cdots$ & Autovalores $0 = \lambda_0 \le \cdots \le \lambda_{n-1}$ \\
\addlinespace
$\hat f(\lambda_k) = \inner{f}{\varphi_k}$ & $\hat f = U^\top f$ \\
\addlinespace
Filtro espectral $\eta(\lambda_k)$ & Filtro espectral $\eta(\Lambda)$ \\
\addlinespace
Elemento de volumen $dV_g$ & Pesos de v\'ertice / \'areas de Voronoi \\
\bottomrule
\end{tabular}
\end{table}

\paragraph{Perspectiva.}
Las geod\'esicas completan el cuadro geom\'etrico iniciado con rejillas, grupos y
grafos: dotan a un dominio de distancia intr\'inseca, curvatura y una noci\'on
rigurosa de flujo de informaci\'on. Dos hilos contin\'uan adelante. La
ambig\"uedad gauge de las coordenadas normales (\Cref{rem:geo:gauge}) --- la
ausencia de un marco local can\'onico en un dominio curvo --- es precisamente el
problema que resuelven las redes equivariantes gauge del \Cref{ch:gauges}. Y el
lenguaje espectral compartido de la \Cref{tab:geo:analogy}, junto con la
agregaci\'on por pseudo-coordenadas de MoNet, sit\'ua las convoluciones en
rejillas, grafos y variedades dentro del \'unico marco de paso de mensajes al
que vuelven los \Cref{ch:applications,ch:conclusion}.
